\section{Setting Up the Forensic Environment}

This laboratory can reuse the CAINE14/Kali VM already configured for a previous forensic exercise. \\
For completeness, and so that this laboratory can be carried out on its own, even if you no longer have that VM available, the full environment setup is summarised again below.

\subsection{Quick VM Setup via Pre-built \texttt{.ova} Image}

If you do not already have a working CAINE14/Kali VM, the fastest path is to import an official
pre-built \texttt{.ova} appliance rather than installing from an ISO. This is the recommended
route if you are starting from scratch for this laboratory.

\begin{enumerate}[leftmargin=1.2cm]
	\item \textbf{Download.} Get the official Kali Linux VirtualBox \texttt{.ova} image from
	      \url{https://www.kali.org/get-kali/#kali-virtual-machines} (choose the VirtualBox
	      variant). Verify the download's checksum against the value published on that same
	      page before proceeding\\
		  \textbf{do not skip this step;} it is the same integrity-verification
	      principle used throughout this course for evidence, applied here to your own tooling.
	\item \textbf{Import.} Open VirtualBox $\rightarrow$ \emph{File} $\rightarrow$ \emph{Import
	      Appliance} $\rightarrow$ select the downloaded \texttt{.ova} file $\rightarrow$
	      \emph{Next}.
	\item \textbf{Adjust resources before finishing the import.} On the settings review screen,
	      increase \textbf{RAM to at least 4096 MB} (the default is often lower), see the
	      warning below for why this matters more here than in previous laboratories.
	\item \textbf{Enable nested virtualisation.} After import, before first boot: select the VM
	      $\rightarrow$ \emph{Settings} $\rightarrow$ \emph{System} $\rightarrow$
	      \emph{Processor} tab $\rightarrow$ tick \textbf{``Enable Nested VT-x/AMD-V''}. This
	      single setting is what allows Docker (Section 3.3) to run correctly inside a VM;
	      without it, Docker containers may fail to start or run at very degraded performance.
	\item \textbf{First boot and login.} Start the VM. Default credentials for the official Kali
	      \texttt{.ova} are \texttt{kali} / \texttt{kali}. \\
		  Change the password immediately after
	      first login (\texttt{passwd}) if this VM will be kept beyond the laboratory session.
	\item \textbf{Confirm networking.} Once logged in, verify the VM has a working network
	      connection before continuing: run \texttt{ip a} (you should see an interface with an
	      assigned IP address, not only \texttt{lo}) and then \texttt{ping -c 3 8.8.8.8} to
	      confirm outbound connectivity.
\end{enumerate}

\begin{warningbox}
If your VirtualBox host does not expose ``Enable Nested VT-x/AMD-V'' at all (greyed out), your
host CPU or hypervisor may not support nested virtualisation. In that case, Docker inside this VM
is not guaranteed to work reliably, and you should flag this to your team.
\end{warningbox}

\subsection{Base Forensic VM (recap)}

If you are instead reusing an existing installation rather than the quick \texttt{.ova} route above, the following recap still applies:

\begin{description}[leftmargin=1.5cm, style=nextline]
	\item[Choice of distribution] CAINE 14 ``Lightstream'' \emph{or} Kali Linux (Forensic
	      Mode). Either is acceptable; both are pre-approved.
	\item[VM configuration] RAM: 4GB minimum (3GB absolute minimum, see the warning below
	      regarding Docker's additional memory footprint).\\
		  Network: host-only or NAT adapter,
	      as long as outbound connectivity works (Docker needs to pull images from the internet
	      the first time each is used).
	\item[Write-blocking] On CAINE14, devices are read-only (\texttt{RO=1}) by default via
	      \texttt{rbfstab}; unlock explicitly with \texttt{blockdev --setrw} only when needed.
	      On Kali, write-blocking is \emph{not} automatic: run
	      \texttt{sudo blockdev --setro /dev/sdX} on any evidence device before contact.
\end{description}

\begin{warningbox}
This laboratory does not involve seized physical media, so write-blocking procedures are not
exercised directly here, but the same VM, once configured, is reused for the Docker/MinIO
setup below. 4GB of RAM should be treated as a firm minimum, not merely a suggestion: Section 3.3
adds a containerised service on top of the base VM's own overhead, and low-RAM VMs have been
observed to make the container unstable or very slow to respond.
\end{warningbox}

\subsection{Installing Docker in the Forensic VM}

Run each of the following lines \textbf{one at a time}, waiting for each to complete before
typing the next, \textbf{do not paste all four lines} as a single block: on some terminal configurations (observed on Kali's default \texttt{zsh}), pasting multi-line blocks can trigger bracketed-paste artefacts that corrupt the input and silently skip commands.

\begin{lstlisting}
sudo apt update
\end{lstlisting}
\begin{lstlisting}
sudo apt install docker.io
\end{lstlisting}
\begin{lstlisting}
sudo systemctl enable --now docker
\end{lstlisting}
\begin{lstlisting}
sudo usermod -aG docker $USER
\end{lstlisting}

\begin{warningbox}
\textbf{Stop here and restart your session.} The \texttt{usermod} command above only takes
effect the \emph{next} time you log in, it does not apply retroactively to your current
shell. Log out and log back in (or reboot the VM) before continuing, otherwise every subsequent
\texttt{docker} command in this section will fail with a permission error unless prefixed with
\texttt{sudo}.
\end{warningbox}

\noindent
After logging back in, confirm Docker actually works before moving on:
\begin{lstlisting}
docker run hello-world
\end{lstlisting}
This should print a short ``Hello from Docker!'' confirmation message. If it hangs, times out, or
errors out, revisit the nested-virtualisation setting from Section 3.1 before proceeding further. This is the most common cause of failure at this exact step.
\medskip

\noindent
Record the Docker version you are running, for the reproducibility reasons set out in Section 2.4:
\begin{lstlisting}
docker --version
\end{lstlisting}

\subsection{Setting Up MinIO}

With Docker confirmed working, start MinIO, that is a local, open-source server that speaks the same S3 API as a real cloud provider, so that the \texttt{aws cli} commands you will run in Section 4 work identically to how they would against a real bucket.
\medskip

\noindent
The image tag below is pinned to the release under which this laboratory was verified, rather than
\texttt{latest}, so that everyone obtains the same behaviour (Section 2.4). The
\texttt{-\/-add-host} flag is needed later, in Section 4.5, so that the container can reach a
listener running on the VM itself:

\begin{lstlisting}
docker run -d --name minio \
  -p 9000:9000 -p 9001:9001 \
  --add-host=host.docker.internal:host-gateway \
  -e "MINIO_ROOT_USER=minioadmin" \
  -e "MINIO_ROOT_PASSWORD=minioadmin" \
  minio/minio:RELEASE.2025-09-07T16-13-09Z server /data --console-address ":9001"
\end{lstlisting}

\noindent
Verify the container started correctly and \emph{stayed running} (this matters: some containers start and then immediately exit due to a configuration error, a running container is not the same as a successfully created one):

\begin{lstlisting}
docker ps
\end{lstlisting}
You should see a line for \texttt{minio} with status \texttt{Up}.\\
If instead \texttt{docker ps} shows nothing, run \texttt{docker ps -a} (which also lists stopped/exited  containers) and then inspect why it exited:

\begin{lstlisting} 
docker logs minio --tail 20
\end{lstlisting}

\noindent
The log output should show the MinIO startup banner (version, API/console addresses) with no
license or authentication error.

\begin{warningbox}
MinIO is a \emph{local, open-source S3-compatible emulator}, not a real cloud environment. It is used here only to give you first-hand experience with the interaction pattern (CLI $\rightarrow$ API $\rightarrow$ object store), not to reproduce the full complexity of a real provider. All log-based analysis in this laboratory relies on artefacts provided separately (Section \ref{sec:log-analysis}), reflecting the fact that a real forensic investigator receives log exports, not raw infrastructure access.
\end{warningbox}

\subsection{Configuring \texttt{aws cli} Against the Local Endpoint}

The \texttt{aws cli} is a generic client: the same tool talks to a real AWS account or to any
S3-compatible endpoint, depending only on which credentials and endpoint URL you give it. Install
it first if not already present (it ships by default on CAINE14; on Kali it must be installed
explicitly):
\begin{lstlisting}
sudo apt install awscli
\end{lstlisting}

\noindent
Record its version too, then configure it with the MinIO root credentials. Run each line separately, exactly as in
Section 3.3, do \emph{not} paste them as one block:
\begin{lstlisting}
aws --version
\end{lstlisting}
\begin{lstlisting}
aws configure set aws_access_key_id minioadmin
\end{lstlisting}
\begin{lstlisting}
aws configure set aws_secret_access_key minioadmin
\end{lstlisting}
\begin{lstlisting}
aws configure set region eu-west-1
\end{lstlisting}

\noindent
Verify the credentials were actually stored before moving on, do not assume the three commands
above succeeded silently: 

\begin{lstlisting}
aws configure list
\end{lstlisting}

\noindent
You should see non-empty (masked) values for \texttt{access\_key} and \texttt{secret\_key}. If either shows as empty, repeat the corresponding \texttt{aws configure set} command above.
\\
Finally, define the alias used throughout the rest of this laboratory:
\begin{lstlisting}
alias awslocal="aws --endpoint-url=http://localhost:9000"
\end{lstlisting}

\begin{warningbox}
This \texttt{alias} only lasts for the current shell session. If you close the terminal or open a
new one, you will need to redefine it (or add it to \texttt{\textasciitilde/.bashrc} /
\texttt{\textasciitilde/.zshrc} if you want it to persist).
\end{warningbox}

\begin{warningbox}
\textbf{An unorthodox practice, flagged as such.} The two commands above put the MinIO root
credentials on the command line in cleartext. This is acceptable \emph{here} only because they are
throwaway credentials for a local, disposable sandbox holding no real evidence. In any real
environment it would be poor practice, and for reasons that matter forensically as much as
operationally: the credentials remain visible in your shell history
(\texttt{\textasciitilde/.zsh\_history}), in the process list while the command runs
(\texttt{ps aux}), and, for the values passed to \texttt{docker run}, in the container metadata
retrievable by any local user with \texttt{docker inspect minio}. An investigator who works this
way on a live system both exposes the credentials and writes avoidable artefacts onto the machine
under examination.\\
Two straightforward alternatives: write the credentials directly into
\texttt{\textasciitilde/.aws/credentials} with restrictive permissions
(\texttt{chmod 600}), so they never transit the command line; or, for the container, place them in
an \texttt{.env} file (also \texttt{chmod 600}) and pass it with
\texttt{docker run -\/-env-file}, which keeps them out of \texttt{docker inspect} output.\\
Keep this in mind for Section 8: the same question, who can see and who can alter the material
your conclusions rest on, applies both to the credentials you use and to the evidence you are
given.
\end{warningbox}

\begin{questionbox}
\texttt{aws cli} does not distinguish between a real AWS account and a local emulator: it is the configuration, endpoint, credentials, region, that decides what it actually talks to, not the tool itself. \\

The credentials you just configured (\texttt{minioadmin}/\texttt{minioadmin}) grant full administrative access to this sandbox bucket. In a real investigation, you would not receive credentials like these: you would be issued access scoped specifically to the case, by the cloud provider, under the terms of the engagement. \\

What does the client's inability to distinguish these two situations tell you about how much trust you can place in a command's output, if you do not independently know how the person who ran it configured their access? And why does the gap between ``administrative access'' and ``access scoped for this investigation'' matter for chain of custody, specifically?
\end{questionbox}

\begin{answerbox}
\vspace{2.5cm}
\end{answerbox}
